\section{Introduction}
\label{sec:intro}

A Planner decomposes a problem, an Executor carries out the plan, and a
Reviewer critiques the answer: this division of labor is a familiar design
for multi-agent large language model systems~\cite{hong2024metagpt,qian2024chatdev,wu2024autogen}.
Yet a role prompt specifies more than the work an agent should do. In a
Planner--Executor--Reviewer (\PER{}) team, the Reviewer can also decide
whether an answer is submitted for external evaluation. This decision
controls an opportunity to learn from a mistake. Our central finding is that
roles can shape failure recovery through these decision rights even when
their intended division of labor provides no first-pass accuracy gain in
the tested settings.

Role instructions do not by themselves establish that agents specialize.
A model may honor the intended division of labor or solve the entire
problem in every role. We therefore measure both the work agents actually
do and the decisions that determine their access to corrective information.
A final accuracy score alone cannot separate these two functions.

The evaluator in our experiments has access to the reference answer and
returns a verdict and, after a failure, a repair hint. Internal deliberation
has no access to this external signal. A conservative Reviewer saves
evaluations, which can be valuable when they are expensive, but also limits
opportunities to obtain corrective feedback. We ask how this trade-off
affects recovery under a fixed call budget.

To isolate the effect of roles, we fix the agent topology, the actor model
(the language model that plays all three agents), the call budget, and the
ground-truth evaluator. We then vary the organization in two directions.
First, we completely delete every role instruction. Second, we strictly
enforce the division of labor at the interface using typed message formats.
Because feedback typically reaches each agent through its prompt, a role
ablation also rewrites the feedback channel. We therefore carefully record
the exact information each condition receives (\S\ref{sec:method}). In our
main comparison, both teams have equal access to the evaluator's verdict
and hint. While access is matched, the realized number and content of
verdicts differ. This difference is what roles control.

Deleting roles largely dismantles the intended division of labor; the
role-free prompt asks each agent for its current answer. Consistent with
prior role ablations~\cite{dylan2026equal,orogat2026architectural}, first-pass
accuracy does not decline. The difference emerges after a rejected first
answer. With feedback available to both teams, the role-free team recovers
more often across all five model--benchmark settings and achieves higher
final accuracy. Submission logs show that the role-structured Reviewer
frequently withholds approval, including before any feedback exists, and
its team consequently receives fewer verdicts. A post hoc comparison at
equal numbers of verdicts yields a small pooled recovery difference
($+0.009$, 95\% interval $[-0.014,+0.033]$). In an exploratory intervention
on $400$ Omni-MATH problems with gemma-4, forcing submission in each cycle
while preserving the role prompts raises recovery from $0.198$ to $0.613$;
the paired comparison detects no difference from the role-free team.
Together, these results support examining submission policy alongside
specialization. They do not establish that the two organizations use each
verdict equally well.

Strictly enforcing the division of labor provides a complementary test.
It loosens the Reviewer gate and raises recovery, with a statistically significant gain only
for Llama-3.3. However, $8\%$ to $17\%$ of enforced trajectories exhaust the
allowed regenerations without finishing, so end-to-end accuracy does not
improve. Across these comparisons, role fidelity and useful feedback access
are distinct properties: a team can follow its assigned roles while losing
opportunities to correct an answer.
We summarize our contributions as follows:
\begin{itemize}
\item[(1)] We introduce a budget-matched design that varies organization in
  two directions and strictly documents the information each condition
  receives. This includes a feedback-withheld diagnostic that shows what a
  role ablation measures when it also removes the feedback channel.
\item[(2)] We provide empirical evidence across three actors that roles
  successfully induce functional specialization, and across two domains
  that they do not improve first-pass accuracy.
\item[(3)] We provide evidence that the Reviewer's approval governs the
  team's feedback opportunities. We support this claim using submission
  logs, a post hoc verdict-matched comparison, and an exploratory
  intervention that removes the submission gate while keeping role prompts
  intact.
\item[(4)] We outline implications for evaluating multi-agent
  organizations. We recommend reporting recovery rates after a failed first
  attempt and treating decision rights and information pathways as core
  components of the organizational structure (\S\ref{sec:discussion}).
\end{itemize}
